\subsection{Ứng dụng \protect\tr{kernel} trong SVM}

\figurename~\ref{fig:kernel-in-svm} minh họa biên quyết định (decision
boundary) của mô hình SVM khi sử dụng bốn loại \tr{kernel} khác nhau gồm:
\tr{kernel} tuyến tính, \tr{kernel} đa thức, \tr{kernel} RBF và \tr{kernel}
sigmoid. Các điểm dữ liệu thuộc hai lớp được biểu diễn bằng hai màu khác nhau,
trong khi đường liền nét thể hiện siêu phẳng phân tách và các đường nét đứt
biểu diễn biên của mô hình.

\begin{itemize}[font=\bfseries,align=left]
    \item[\tr{Kernel} tuyến tính:] Ranh giới quyết định là một đường thẳng, phù
    hợp với dữ liệu có khả năng phân tách tuyến tính. Tuy nhiên, với tập dữ
    liệu có sự phân bố phức tạp, \tr{kernel} tuyến tính chưa thể phân loại tối
    ưu.

    \item[\tr{Kernel} đa thức:] Ranh giới quyết định trở nên cong
    hơn, cho phép mô hình học được các quan hệ phi tuyến giữa các lớp.
    \tr{Kernel} này cải thiện khả năng phân loại so với \tr{kernel} linear
    nhưng vẫn còn hạn chế ở những vùng dữ liệu chồng lấn mạnh.

    \item[\tr{Kernel} RBF:] Ranh giới quyết định linh hoạt và phi tuyến mạnh
    mẽ nhất. Mô hình có khả năng bao quát tốt cấu trúc dữ liệu phức tạp, giúp
    phân tách hai lớp hiệu quả hơn. Đây thường là \tr{kernel} cho kết quả tốt
    rong nhiều bài toán thực tế.

    \item[\tr{Kernel} Sigmoid:] Ranh giới quyết định có dạng tương tự mạng
    nơ-ron với các vùng phân tách phức tạp. Tuy nhiên, kết quả phân loại trong
    hình cho thấy \tr{kernel} sigmoid không ổn định bằng RBF và có thể tạo ra
    các biên phân tách chưa tối ưu.
\end{itemize}

\begin{figure}[htbp]
  \centering
  \includegraphics[width=0.9\textwidth]{kernel-in-svm}
  \caption{Đường biên phân loại SVM trong các không gian đặc trưng}
  \label{fig:kernel-in-svm}
\end{figure}

Từ kết quả trực quan có thể thấy rằng việc lựa chọn \tr{kernel} ảnh hưởng lớn
đến hiệu suất phân loại của SVM. Với dữ liệu phi tuyến, các \tr{kernel} như
\tr{kernel} đa thức và đặc biệt là RBF thường cho khả năng phân tách tốt hơn
so với \tr{kernel} tuyến tính.
